% ============================================================================
% Chapter 1 --- Synthesis (modular file)
% ----------------------------------------------------------------------------
% Purpose : a self-contained ~12-page synthesis of the entire Technical
%           Reference. Read on its own when this chapter is included as the
%           book's Part IV; serves as executive overview when read at the
%           head of the standalone Technical Reference.
% Audience: CTO, RSSI, governance, anyone who needs the essentials in 30 min.
% Page target: ~12 pages.
% Cross-refs use the \trchap macro so phrasing stays unambiguous in both
% inclusion contexts (standalone and book).
% ============================================================================

\chapter{Synthesis}
\label{ch:synthesis}

This chapter is a self-contained executive overview of the entire
Technical Reference. It develops, at the level of essential
content rather than implementation detail, the three structural
contributions of the document: the sovereignty grid, the
reference architecture on open foundations, and the wave-based
migration playbook. It also provides a one-minute
self-assessment, an annotated set of institutional precedents,
and a reading map for the remaining thirty-two chapters. A
reader who completes this chapter alone should be able to decide
which subsequent chapters are most relevant to their role.

%-----------------------------------------------------------------------------
\section{Purpose and audience}
\label{sec:synth-purpose}

The Technical Reference addresses chief technology officers,
information-security officers, network and platform engineers,
and the institutional governance bodies that mandate them. It
complements the policy-and-analysis main report with the
material a practitioner would need to operationalise the
architecture: explicit interface contracts, a vendor-neutral
reference instantiation on open foundations, alternatives that
admit proprietary components within the same contracts, and a
migration playbook applicable across institutions and
jurisdictions.

\paragraph{Reading time and dependence.} The present chapter is
designed to be read in approximately thirty minutes by a reader
familiar with university IT. It depends on no prior knowledge
of the Reference's other chapters; cross-references throughout
point the reader to the development that they may wish to
consult next.

\paragraph{Where this chapter sits in the suite.} The Technical
Reference is one of the four documents of the suite on university
IT infrastructure: the policy dossier --- a policy-and-analysis
main report with its executive-summary companion and its
question-and-answer part --- and the present
practitioner-oriented Technical Reference, with its operational
summary and its own question-and-answer module. The dossier does
not reproduce this chapter: it carries a pointer-and-capsule part
orienting its readers to the standalone Reference, of which this
chapter is the opening chapter.

%-----------------------------------------------------------------------------
\section{The TAC concept at a glance}
\label{sec:synth-tac}

A \keyword{Trusted Activity Context} (\gls{tac}) is a coherent,
validated environment that assembles the data, tools, services,
and access rights required for a specific activity --- a research
project, a digital examination, an administrative process, a
teaching session --- under documented institutional governance.
Within a TAC, the identity of subjects, the classification of
resources, the network attachment, the policy decisions, the
cryptographic posture, and the audit trail are aligned with the
activity's objectives and risk level, and are operationally
bounded by the TAC's declared scope.

\paragraph{Why it matters.} Existing access-control and isolation
frameworks are typically organised around either subjects (who a
user is) or resources (what a service contains). A TAC is
organised around \emph{what is being done}: the same subject,
holding the same attributes, accessing the same resource, may
operate in a different TAC depending on the activity engaged.
This re-frames several recurring difficulties in academic IT ---
the researcher who is also a teacher, the visiting collaborator
who is also a host, the administrative process that touches
research data --- as switches between TACs rather than as
exceptions to a single role-based model.

\begin{figure}[H]
\centering
\begin{tikzpicture}[
    every node/.style = {font=\footnotesize},
    role/.style = {rectangle, rounded corners=2pt,
                   draw=NavyBlue!70!black, fill=NavyBlue!5,
                   minimum width=18mm, minimum height=8mm,
                   line width=0.4pt, align=center},
    arrow/.style = {-{Stealth[length=1.8mm]}, draw=NavyBlue!50!black,
                    line width=0.4pt}
]
\node[role] (subject) at (0,0)   {Subject};
\node[role] (id)      at (2.2,0) {Identity};
\node[role] (po)      at (4.4,0) {Policy};
\node[role] (nw)      at (6.6,0) {Network};
\node[role] (co)      at (8.8,0) {Compute};
\node[role] (au)      at (11.0,0) {Audit};

\draw[arrow] (subject) -- (id);
\draw[arrow] (id) -- (po);
\draw[arrow] (po) -- (nw);
\draw[arrow] (nw) -- (co);
\draw[arrow] (co) -- (au);

% Correlation_id thread
\node[font=\scriptsize, anchor=west] at (-0.5,-1.4) {%
  TAC \texttt{correlation\_id} propagated across all steps:};
\draw[-, draw=Crimson!70!black, line width=1pt, dashed]
  (-0.5,-1.7) -- (12.0,-1.7);
\foreach \x in {0,2.2,4.4,6.6,8.8,11.0} {
  \draw[draw=Crimson!70!black, line width=0.5pt] (\x,-0.5) -- (\x,-1.7);
}
\end{tikzpicture}
\caption{TAC activation sequence (synthesis view): the activation
chain crosses six functional roles, with the
\texttt{correlation\_id} threading every step for cross-layer
audit reconstruction. Full eleven-step development in
\trchap{ch:conceptual-frame}{} (\S\ref{sec:cf-activation-sequence}).}
\label{fig:synth-tac-sequence}
\end{figure}

\paragraph{Pointer.} The full conceptual development of TAC ---
formal definition, principles, comparison with related
frameworks, lifecycle state machine --- is in
\trchap{ch:conceptual-frame}{}.

%-----------------------------------------------------------------------------
\section{The sovereignty grid in one table}
\label{sec:synth-sovereignty}

Sovereignty in university IT is rarely contested in principle and
rarely operationalised in practice. The grid that the Technical
Reference proposes is the instrument that bridges the two: a
structured way of asking, for any candidate component of the
digital infrastructure, whether the institution retains the
capacities --- legal, technical, financial, operational, and
institutional --- that allow it to remain the author of its own
decisions.

\paragraph{Five dimensions, each scored independently.} The grid
is structurally non-aggregating: each dimension is reported on
its own four-point scale (Absent / Partial / Established /
Robust); no single sovereignty score is produced. A single
critical gap on a dimension defines the institution's exposure on
that dimension regardless of the other dimensions' strength.

\begin{table}[H]
\centering\small
\caption{The five dimensions of the sovereignty grid (synthesis;
full version with twenty indicators in
\trchap{ch:sovereignty-grid}{}).}
\label{tab:synth-sovereignty}
\begin{tabularx}{\textwidth}{%
  >{\hsize=0.4\hsize\linewidth=\hsize\bfseries\raggedright\arraybackslash}X%
  >{\hsize=1.6\hsize\linewidth=\hsize\raggedright\arraybackslash}X%
}
\toprule
Dimension & Question (verifiable capacity) \\
\midrule
Data
  & Where is the data stored, under which jurisdiction, accessible
    to which actors --- including state actors? \\
\addlinespace
Technical
  & Can the institution inspect, modify, repair, and substitute
    the components of its infrastructure? \\
\addlinespace
Financial
  & Is the multi-year operating cost predictable, and is the
    institution exposed to captive escalation? \\
\addlinespace
Operational
  & Are the skills required to operate the infrastructure
    available locally or through trusted partners (the national
    research and education network, peer institutions)? \\
\addlinespace
Institutional
  & Can the institution decline a supplier roadmap change without
    contractual penalty, and retain authority over the decisions
    that matter? \\
\bottomrule
\end{tabularx}
\end{table}

\paragraph{Use as a governance instrument.} The grid feeds four
governance practices that are frequently
conducted in mutual ignorance: procurement (RFP language,
contractual exit assistance, gating thresholds per workload
class); internal audit (annual review of the institution's
sovereignty profile per component); peer-institution benchmarking
(comparison anchored by a shared grid); and public accountability
(publication of the institution's sovereignty position by
component).

\paragraph{Pointer.} The grid is fully developed in
\trchap{ch:sovereignty-grid}{} (twenty indicators, scoring
guidance, three worked examples on identity instantiations,
common misuses, governance use). The institutional self-assessment
that operationalises the grid is in
\trchap{ch:self-assessment}{}.

%-----------------------------------------------------------------------------
\section{Reference architecture in one table}
\label{sec:synth-architecture}

The Technical Reference proposes a vendor-neutral architecture
organised in nine functional layers grouped in three planes ---
Subject and Trust (Identity, Endpoint, Crypto), Execution
(Network, Compute, Storage), and Control and Governance (Policy,
Observability, Orchestration). The architecture's normative
content is the set of interface contracts (open standards) that
connect the layers; products are illustrative, not prescriptive.

\paragraph{The reference instantiation on open foundations.} A
reference instantiation is given on open-source components to
demonstrate that the contracts are realisable without proprietary
extensions in critical paths. Commercial alternatives that
respect the contracts remain admissible; the choice between
instantiations is informed by the sovereignty grid rather than
imposed.

\begin{figure}[H]
\centering
\begin{tikzpicture}[
    every node/.style = {font=\footnotesize},
    layer/.style = {rectangle, rounded corners=2pt, draw=NavyBlue!70!black,
                    fill=NavyBlue!5, minimum width=22mm, minimum height=8mm,
                    line width=0.4pt, align=center},
    plane/.style = {rectangle, rounded corners=4pt, draw=Slate!50,
                    line width=0.4pt, dashed, inner sep=4pt},
    arrow/.style = {-{Stealth[length=1.8mm]}, draw=NavyBlue!50!black,
                    line width=0.3pt}
]
% Subject and Trust plane (top-left)
\node[layer] (id)   at ( 0.0, 0.0) {Identity};
\node[layer] (ep)   at ( 2.6, 0.0) {Endpoint};
\node[layer] (cr)   at ( 5.2, 0.0) {Crypto};
\node[plane, fit=(id)(ep)(cr), label={[font=\scriptsize\bfseries,
      color=Slate!80!black]above:Subject and Trust plane}] {};

% Execution plane (middle)
\node[layer] (nw)   at ( 0.0,-2.6) {Network};
\node[layer] (co)   at ( 2.6,-2.6) {Compute};
\node[layer] (st)   at ( 5.2,-2.6) {Storage};
\node[plane, fit=(nw)(co)(st), label={[font=\scriptsize\bfseries,
      color=Slate!80!black]left:Execution plane}] {};

% Control and Governance plane (bottom)
\node[layer] (po)   at ( 0.0,-5.2) {Policy};
\node[layer] (ob)   at ( 2.6,-5.2) {Observ.};
\node[layer] (or_)  at ( 5.2,-5.2) {Orchestr.};
\node[plane, fit=(po)(ob)(or_), label={[font=\scriptsize\bfseries,
      color=Slate!80!black]below:Control and Governance plane}] {};

% Principal flows
% Subject/Trust -> Control (Identity attributes feed Policy)
\draw[arrow] (id.south) to[bend right=15] (po.north);
\draw[arrow] (ep.south) to[bend right=10] (po.north);
% Control -> Execution (Policy decisions enforced at Network/Compute/Storage)
\draw[arrow] (po.east) -- (nw.south);
\draw[arrow] (po.east) to[bend right=10] (co.south);
\draw[arrow] (po.east) to[bend right=18] (st.south);
% All -> Observability
\draw[arrow] (nw.south east) to[bend right=8] (ob.north);
\draw[arrow] (co.south) -- (ob.north);
\draw[arrow] (st.south) to[bend left=8] (ob.north east);
% Orchestration spans
\draw[arrow] (or_.north) to[bend left=15] (st.south east);
\end{tikzpicture}
\caption{Layer-dependency map of the reference architecture: nine
functional layers in three planes. Subject and Trust feeds the
Control plane; the Control plane writes to the Execution plane;
Observability collects events from every layer; Orchestration
spans the Execution plane on activation events. Full development
in \trchap{ch:architectural-principles}{} through
\trchap{ch:orchestration}{}.}
\label{fig:synth-arch-layers}
\end{figure}

\begin{table}[H]
\centering\small
\caption{Reference architecture summary: nine functional layers
with their principal interface contract and reference open
instantiation (synthesis; per-layer detail in
Chapters~6--14 of the Technical Reference).}
\label{tab:synth-architecture}
\begin{tabularx}{\textwidth}{%
  >{\hsize=0.45\hsize\linewidth=\hsize\bfseries\raggedright\arraybackslash}X%
  >{\hsize=0.95\hsize\linewidth=\hsize\raggedright\arraybackslash}X%
  >{\hsize=1.60\hsize\linewidth=\hsize\raggedright\arraybackslash}X%
}
\toprule
Layer & Interface contract (standards) & Open reference \\
\midrule
Identity        & OIDC, SAML2, SCIM 2.0, federation metadata  & Keycloak + SCIM + eduGAIN \\
Policy          & OPA REST or Cedar evaluation API            & OPA + bundle server + decision log \\
Network         & RADIUS, NETCONF/YANG, BGP/EVPN, IPFIX       & VLAN/VRF + FreeRADIUS + Open vSwitch + research-data path (Science DMZ / national research and education network circuit) \\
Compute         & Kubernetes API, OpenStack API, OCI, SLSA    & OpenStack + Kubernetes + Harbor + SLSA \\
Storage         & S3 API, KMIP, NFS v4.1                      & Ceph + Vault + label-based access \\
Endpoint        & DDM/DEP, MS-MDM, declarative Linux config   & Munki + osquery + Ansible + MicroMDM (gap on iOS/Android at scale acknowledged) \\
Observability   & OpenTelemetry, CloudEvents, syslog          & Vector + Kafka + OpenSearch + Grafana Loki + Grafana \\
Crypto          & PKCS\#11, ACME, KMIP                        & step-ca + Vault + SoftHSM or Nitrokey HSM \\
Orchestration   & Event-driven, idempotent, observable, reversible (no single standard) & Argo Workflows + Crossplane + GitOps \\
\bottomrule
\end{tabularx}
\end{table}

\paragraph{Conformance and replaceability.} The architecture's
contracts are operationalised by a conformance suite
organised in four categories (interface
fidelity, data portability, exit rehearsal, decision auditability).
Every wave of the migration playbook closes by passing the
conformance subset relevant to its scope. The suite is the
mechanism that prevents the common failure mode of migration
programmes: change of product disguised as architectural progress.
Full description in \trchap{ch:conformance}{}.

%-----------------------------------------------------------------------------
\section{Migration playbook in one figure and one table}
\label{sec:synth-migration}

The Technical Reference proposes a wave-based migration playbook
that takes any institution from its current brownfield position
toward the target architecture. Ten waves (Wave~0 to Wave~9)
unfold under a documented dependency graph. Each wave declares
its goal, pre-requisites, coexistence pattern, three principal
trajectories (greenfield, brownfield single-suite-centric --- a
dominant proprietary productivity suite --- and brownfield hybrid),
reversibility criteria,
and exit criteria gated by the conformance suite.

\paragraph{Three pace options.} The playbook admits three reference
pace options --- 5-year compressed, 7-year balanced (recommended
default), 10-year gradual. Pace and trajectory interact; revising
the pace is a normal part of operating the playbook.

\begin{figure}[H]
\centering
\begin{tikzpicture}[
    every node/.style = {font=\footnotesize},
    wave/.style = {circle, draw=NavyBlue!70!black, fill=NavyBlue!5,
                   minimum size=11mm, line width=0.5pt, align=center,
                   inner sep=1pt},
    arrow/.style = {-{Stealth[length=1.8mm]}, draw=NavyBlue!50!black,
                    line width=0.3pt}
]
% Layer 0: Foundations
\node[wave] (w0) at (0,0) {0};
% Layer 1: parallel waves 1, 2, 3
\node[wave] (w1) at (-3.0,-1.6) {1};
\node[wave] (w2) at (-1.0,-1.6) {2};
\node[wave] (w3) at ( 1.0,-1.6) {3};
% Layer 2: wave 4 (needs 1+2)
\node[wave] (w4) at (-2.0,-3.2) {4};
% Layer 3: wave 5 (pilot, needs 1+2+3+4) and waves 6, 7, 8
\node[wave] (w5) at (-2.0,-4.8) {5};
\node[wave] (w6) at ( 0.5,-3.2) {6};
\node[wave] (w7) at ( 0.5,-4.8) {7};
\node[wave] (w8) at ( 3.0,-3.2) {8};
% Layer 4: wave 9 (closes on cumulative)
\node[wave, fill=NavyBlue!15] (w9) at ( 3.0,-4.8) {9};

% Wave 0 to 1, 2, 3
\draw[arrow] (w0) -- (w1);
\draw[arrow] (w0) -- (w2);
\draw[arrow] (w0) -- (w3);
% Wave 1+2 to 4
\draw[arrow] (w1) -- (w4);
\draw[arrow] (w2) -- (w4);
% Wave 1+2+3+4 to 5
\draw[arrow] (w1) to[bend right=15] (w5);
\draw[arrow] (w2) to[bend right=10] (w5);
\draw[arrow] (w3) -- (w5);
\draw[arrow] (w4) -- (w5);
% Wave 1+4 to 6
\draw[arrow] (w1) to[bend left=10] (w6.west);
\draw[arrow] (w4) -- (w6);
% Wave 2+4+5 to 7
\draw[arrow] (w2) to[bend left=10] (w7);
\draw[arrow] (w4) to[bend right=10] (w7);
\draw[arrow] (w5) -- (w7);
% Wave 1+3+4 to 8
\draw[arrow] (w1) to[bend left=20] (w8);
\draw[arrow] (w3) to[bend left=10] (w8);
\draw[arrow] (w4) to[bend left=15] (w8);
% Wave 1-8 to 9
\draw[arrow] (w6) -- (w9);
\draw[arrow] (w7) -- (w9);
\draw[arrow] (w8) -- (w9);

% Labels
\node[font=\scriptsize, anchor=west] at (4.0, 0) {Wave~0 — Foundations};
\node[font=\scriptsize, anchor=west] at (4.0,-1.6) {1 — Identity\quad 2 — Observability\quad 3 — Network};
\node[font=\scriptsize, anchor=west] at (4.0,-3.2) {4 — Policy\quad 6 — Endpoint\quad 8 — Workload};
\node[font=\scriptsize, anchor=west] at (4.0,-4.8) {5 — Pilot\quad 7 — Audit/IR\quad 9 — Lifecycle};
\end{tikzpicture}
\caption{Wave-dependency map of the migration playbook. Wave~0
is the entry point; Waves~1, 2 and 3 can proceed in parallel;
later waves close on cumulative pre-requisites. Full development
in \trchap{ch:playbook-reading}{}.}
\label{fig:synth-wave-dag}
\end{figure}

\begin{table}[H]
\centering\small
\caption{Migration playbook summary: ten waves with goal and
principal exit criterion (synthesis; per-wave detail in
Chapters~18--27 of the Technical Reference).}
\label{tab:synth-waves}
\begin{tabularx}{\textwidth}{%
  >{\hsize=0.75\hsize\linewidth=\hsize\bfseries\raggedright\arraybackslash}X%
  >{\hsize=1.20\hsize\linewidth=\hsize\raggedright\arraybackslash}X%
  >{\hsize=1.05\hsize\linewidth=\hsize\raggedright\arraybackslash}X%
}
\toprule
Wave & Goal & Principal exit criterion \\
\midrule
0 — Foundations            & Governance, baseline inventory, sovereignty grid adoption.
                           & Four artefacts in writing; programme funded. \\
\addlinespace
1 — Identity               & SSO, federation, SCIM aligned with the contract.
                           & Identity conformance subset passes; agreed application coverage. \\
\addlinespace
2 — Observability          & Pre-instrumentation in shadow mode.
                           & Source coverage; canonical schema enforced. \\
\addlinespace
3 — Network                & Zone catalogue, RADIUS attachment, high-throughput research-data path (Science DMZ or national research and education network circuit).
                           & Network conformance subset passes. \\
\addlinespace
4 — Policy plane           & PDP/PEP introduced in shadow before enforcing.
                           & Policy conformance subset passes; migrated set enforcing. \\
\addlinespace
5 — Pilot                  & Full stack exercised end-to-end on a contained scope.
                           & Union of relevant conformance subsets; pilot host report. \\
\addlinespace
6 — Endpoint               & Posture signals routed to policy; per-population transition.
                           & Endpoint conformance subset passes; agreed populations enforcing. \\
\addlinespace
7 — Audit and IR           & TAC-correlated investigation; refreshed runbooks.
                           & Audit conformance subset; SOC proficiency demonstrated. \\
\addlinespace
8 — Workload classification& Classification enforced at provisioning; cluster separation.
                           & Compute and Storage conformance subsets; classification coverage. \\
\addlinespace
9 — Lifecycle automation   & Self-service, GitOps, mature operational stage.
                           & Orchestration conformance subset; self-service ratio achieved. \\
\bottomrule
\end{tabularx}
\end{table}

\paragraph{Reversibility as non-negotiable.} The playbook commits
to the reversibility of every wave at acceptable cost during a
published window. The reversibility playbook of
\trchap{ch:reversibility}{} consolidates the per-wave rollback
procedures, the data-preservation considerations, the
communication patterns, and the post-mortem framework. The
commitment makes the migration a governance decision rather than
a leap of faith.

%-----------------------------------------------------------------------------
\section{Self-assessment in one minute}
\label{sec:synth-self-assessment}

The full institutional self-assessment of
\trchap{ch:self-assessment}{} comprises approximately fifty
questions organised in five sections. The five questions below
are a one-minute version: one question per sovereignty dimension,
each answered \textbf{Y} (capacity established or robust) or
\textbf{N} (absent or partial).

\begin{enumerate}[leftmargin=2em,nosep]
  \item \textbf{Data.} Are storage location, key custody, and
    disclosure regime documented per data class?
  \item \textbf{Technical.} Can the institution inspect or
    independently rebuild every critical component of its
    infrastructure?
  \item \textbf{Financial.} Is the multi-year operating cost of
    the digital infrastructure predictable and contractually
    capped?
  \item \textbf{Operational.} Are the skills for each critical
    layer present in the institution or accessible through its
    national research and education network and peer network?
  \item \textbf{Institutional.} Can the institution decline a
    supplier roadmap change without contractual penalty for the
    current procurement period?
\end{enumerate}

\paragraph{Reading the result.} Five Y answers indicate that the
institution is already operating close to the architecture's
sovereignty target; the playbook is then used as a checklist
rather than as a roadmap. Mixed answers indicate the dominant
trajectory: predominantly N answers on Data, Technical, and
Institutional with a Y on Operational suggest a brownfield
profile dominated by a single proprietary productivity suite;
mixed answers across all dimensions
suggest a brownfield hybrid profile; predominantly N answers
across the board suggest that Wave~0 (governance and inventory)
is the appropriate first step before trajectory selection.

%-----------------------------------------------------------------------------
\section{Real-world precedents at a glance}
\label{sec:synth-precedents}

The architecture and the playbook do not invent their structural
patterns from first principles; they consolidate patterns visible
at named institutions whose practices have been the subject of
public discussion. The list below records the institutions that
have appeared most consistently as references; the full annotated
catalogue, with the methodological caveats, is in
\trchap{ch:institutional-precedents}{}.

\begin{itemize}[leftmargin=1.5em,nosep]
  \item \institution{CERN} --- federated identity at scale; experimental network isolation patterns.
  \item \institution{MIT} --- Information Technology Governance Committee (ITGC); federated IT pattern across institutional and college units.
  \item \institution{Stanford} --- SUNet multi-zone segmentation; UIT and SRCC.
  \item \institution{ETH Z\"urich} --- a binding administrator-rights agreement (a contractual, proportionate transfer of root-access responsibility, on the ETH Z\"urich model); federated IT-services organisation.
  \item \institution{Cambridge} --- University Information Services (UIS); collegiate model with the University Data Network (UDN) connecting central services and the institutions of the Collegiate University.
  \item \institution{Indiana University} --- Jetstream2 research-cloud architecture (NSF ACCESS).
  \item \institution{NUS Singapore} --- Asia-Pacific federation patterns.
  \item \institution{USP Brazil} --- Latin American research-network integration through RNP.
  \item \institution{University of Michigan} --- ITS, Merit Network, CAEN faculty IT.
  \item \institution{KU Leuven} --- European research-intensive central IT model.
\end{itemize}

\paragraph{Documented gaps.} Several aspects of the architecture
remain thinly documented in public institutional sources: policy
plane deployments at academic scale; TAC-aware audit pipelines;
sovereignty grid practice as a structured five-dimension
instrument applied across an institutional portfolio. The
cross-institutional infrastructure-sovereignty programme of
participating institutions anticipated by this Reference is the
locus where these gaps are intended to be closed.

%-----------------------------------------------------------------------------
\section{How to use the rest of the Technical Reference}
\label{sec:synth-using}

The remaining thirty-two chapters of the Technical Reference
develop the elements summarised above. Three reading paths cover
the principal use cases.

\paragraph{Reading path: governance and rectorate.}
Chapter~\ref{ch:conceptual-frame} (the TAC concept),
\trchap{ch:sovereignty-grid}{} (the sovereignty grid in full),
\trchap{ch:self-assessment}{} (the institutional self-assessment).
Indicative reading time: approximately ninety minutes. Output:
the institution's positioning against the grid and an indication
of the dominant migration trajectory.

\paragraph{Reading path: chief technology officer or
information-security officer.} Add to the previous path: the
architectural principles of \trchap{ch:architectural-principles}{};
the layer chapter most relevant to the institution's current
priority among Chapters~6--14; the conformance suite of
\trchap{ch:conformance}{}; and the playbook reading guide of
\trchap{ch:playbook-reading}{}. Indicative reading time:
approximately a working day. Output: a position from which the
institution can engage its trajectory selection and procurement
adjustments.

\paragraph{Reading path: programme or migration lead.} Add to the
CTO path: Part~III in full
(\trchap{ch:playbook-reading}{} through
\trchap{ch:coexistence-patterns}{}). Indicative reading time:
two to three working days. Output: a wave plan with declared
trajectories, paces, and reversibility windows.

\paragraph{Layer-specific reading paths.} An engineer engaged
on a specific layer can read the corresponding chapter (Identity:
\trchap{ch:identity}{}; Policy: \trchap{ch:policy}{}; Network:
\trchap{ch:network}{}; Compute: \trchap{ch:compute}{}; Storage:
\trchap{ch:storage}{}; Endpoint: \trchap{ch:endpoint}{};
Observability: \trchap{ch:observability}{}; Crypto:
\trchap{ch:crypto}{}; Orchestration: \trchap{ch:orchestration}{}),
followed by the corresponding wave chapter. Each layer chapter
is self-contained.

\medskip

\noindent \textbf{End of synthesis.} The remainder of the Technical
Reference (Parts~I to IV) develops each of the elements summarised
above in the depth required for technical implementation,
operational planning, and institutional governance.
